% =====================================================================
% overleaf/chapters/01_Introduction.tex — LaTeX rendering of en/01_Introduction.md v0.27 (22 September 2026)
% Dernière mise à jour : 2026-09-22
% Chapter 1 of the paper. The related work is a \section of its own: overleaf/chapters/02_Related_work.tex,
% to be inserted right after this file. Cross-references use \label / \ref, so no number is hard-coded.
% Ready to paste into the Overleaf project (acmart + natbib: \citep, \citet).
% Derived file: regenerate after every validated pass of en/01_Introduction.md.
% Citation keys are those of docs/paper/sources/sample.bib (checklist: article/CITATIONS.md).
% Preamble additions:
%   \newcommand{\todo}[1]{\textbf{[#1]}}   % placeholder figures, filled from methode/experience_plan/experiments.yaml
% Vocabulary: minimal prompt (Level 1) / expert prompt (Level 2), aligned on the French and on the abstract
%   since v0.25. Level 0--3 = our ablation levels; exploratory / robust / transferable = SILICA certification steps.
% Traceability: the "source:" comments live in the ../fr/ and ../en/ masters only. Do not carry them
%   over into this file when rendering a chapter into LaTeX (author's decision, 2026-09-21).
% =====================================================================

\section{Introduction}
\label{sec:intro}

\subsection{Individual mobility behaviour}

Understanding individual mobility behaviour in multimodal transport systems is essential to travel planning and network design. The challenge is not only to reproduce aggregate flows: it is to build a simulation framework able to support the study of individual behaviour, the testing of fine-grained behavioural hypotheses and the design of personalised mobility policies. For this, the simulated population must reproduce, at least in aggregate, the behaviour recorded by certified household travel surveys \citep{tisseo2023emc2}.

The rise of personalised multimodal transport reflects a transition toward user-centric mobility, in which public transit, walking, cycling and shared services are combined into journeys tailored to individual preferences \citep{smith1995transims,horni2016matsim}. Representing such systems in simulation is difficult because mobility choices are shaped by heterogeneous socioeconomic profiles, accessibility constraints and lived experience \citep{grignard2018impact}. As \citet{oberoi2024personalisation} underlines, personalisation within Mobility-as-a-Service requires dealing with the heterogeneity of users, their preferences, their life contexts and their accessibility needs, which remains a major obstacle to realistic transport models.

Current transport simulations suffer from a shortage of calibration data: agent-based models require extensive local data, which hinders their adoption \citep{liu2024toward}. Detailed multimodal trajectory datasets are scarce and exist only for a few large metropolitan areas, Tokyo for instance \citep{feng2024agentmove,fourez2025transport}. Encoding behavioural rules by hand produces rigid models, costly to scale and helpless in unforeseen conditions \citep{liu2024toward,alves2026evaluating}.

Over the last three years, large language models (LLMs) have entered agent-based simulation as decision engines. Generative agents endowed with natural-language reasoning, episodic memory and planning were introduced by \citet{park2023generative}. They have since been proposed for travel demand modelling \citep{liu2024toward}, for large-scale urban simulation \citep{bougie2025citysim}, and for route replanning under road disruptions on the GAMA platform \citep{alves2026evaluating}. This work builds on the GAMA--OpenTripPlanner--LLM architecture of \citet{vu2025modeling}, applied to the Toulouse metropolitan area.

The appeal of these agents answers the two difficulties above. They promise to represent decisions that no tabulated variable captures: a strike, a heatwave, a news item about safety in the metro. And they promise to do so without the calibration data that classical approaches require: a model that reasons from a persona and a set of itineraries seems to bypass both the scarcity of trajectories and the brittleness of hand-written rules.

Whether it does is an empirical question, and the published answers are partial. Most evaluations compare LLM agents with rule-based agents on scenario plausibility, or on aggregate patterns of their own making. One recent system, GTA \citep{lammer2026gta}, does confront a simulated modal split with a national travel survey, but gives itself no reference beyond the observed split of other regions, whose price Section~\ref{sec:related} shows. What is missing is a protocol under which a measurement becomes attributable: a floor beneath the agents, calibrated tabular models above them, and the same inputs given to every model compared. This paper supplies that protocol and applies it against a certified survey.

\subsection{A population plausible one agent at a time can be wrong in aggregate}
\label{sec:gap}

An LLM agent produces fluent, individually reasonable justifications for its mode choices. A population of such agents may be plausible one agent at a time and wrong in aggregate, because each agent draws on the same pre-training prior rather than on the behaviour of the territory it is meant to represent.

SILICA \citep{bintareaf2026silica} gives this concern an operational form on game-theoretic environments with published human anchors. Twelve open-weight models are run through five games; no finding matches its human anchor by an equivalence test, and convention formation alone survives the perturbations that leave the rules unchanged.
The author bounds the scope: small open-weight models on a consumer graphics card, aggregate anchors obtained mostly on Western university samples, a single-author analysis. Mode choice differs from a repeated game. It is an individual decision conditioned on distance, schedules and vehicle availability, and its reference is a certified survey of the same territory. Whether the ceiling SILICA observes appears on mode choice is an open question.

The alignment literature of Section~\ref{sec:alignment} shows that supervision moves distributions. We read the survey's published modal shares, compare them with the split our agents produce, and rewrite the prompt by hand to close the gap; the prompt carries behavioural guidance in words, never a figure from the survey, a target share or an observed choice. The question is whether a city's own survey is enough to adapt an uncalibrated simulation to that city, or whether the gap is structural. Section~\ref{sec:contributions} phrases it as a hypothesis.

A population of LLM agents can also be asked two different things:

\begin{enumerate}
  \item to serve as a \textbf{statistical stand-in} for the population of a territory, reproducing its modal split across demographic strata;
  \item to serve as an \textbf{adaptive system} that reacts to contexts no survey tabulates, a strike, a heatwave, a service breakdown, and carries the memory of a bad experience from one day to the next.
\end{enumerate}

The first use is judged by distributional fidelity. The second is judged by the robustness and direction of the reaction, because no survey follows the same individuals through such events. We evaluate them separately.

\subsection{Contributions and hypotheses}
\label{sec:contributions}

The study is conducted on the Toulouse metropolitan area, against the 2023 Cerema-certified household travel survey (EMC\textsuperscript{2} 2023) \citep{tisseo2023emc2}, about 16,000 respondents over 453 municipalities.

\paragraph{C1. A comparison protocol on a common set of inputs.}
Three rules, stated in Section~\ref{sec:contract}, each close one way in which a comparison between a generative agent and a tabular model can be biased: the same 21 input variables on both sides, probability mass as the scored quantity rather than a single mode, and the restriction of every tabular prediction to the modes actually offered for the trip. What is equalised is the inputs, not the exposure: the tabular references are fitted on the survey microdata, of which the agent has read no record. At Level~2, the prompt carries behavioural guidance calibrated by hand against the survey's published aggregate shares, never against an individual record.

The evaluation substrate is a sealed synthetic cohort of 1,000 personas
in 499 entire households, the household being the drawing unit, drawn from an eqasim pool of 11,329 persons \citep{horl2021eqasim}. The cohort is controlled against the thirteen margins for which the EMC\textsuperscript{2} survey provides a target, from age band and occupation to transport pass and share of immobile persons, all conformant within a $\pm$\,1 point equivalence bound.
Their activity chains are cyclic, the first ``home'' activity beginning the evening before and the last trip returning home, so that a chain of $k$ activities carries $k$ trips and not $k-1$, and produce 3,299 trips on the evaluated day, i.e.\ 3.30 per persona.
The tabular models are additionally characterised on the sealed survey test set of 13,045 trips, split by household.

\paragraph{C2. An ablation in four levels.}
Four levels of increasing information are compared on the same sealed substrate:

\begin{itemize}
  \item \emph{Level~0}: floors and physical heuristics, a uniform draw over the offer, the majority mode, the shortest itinerary;
  \item \emph{Level~1}: the minimal prompt, factual and circumstantial, with no arbitration instruction;
  \item \emph{Level~2}: the same LLM with the expert prompt of Section~\ref{sec:gap}, obtained by prompt engineering on a calibration population drawn separately from the sealed cohort;
  \item \emph{Level~3}: the four tabular references fitted on the survey (Section~\ref{sec:references}).
\end{itemize}

Several LLMs are evaluated, all at temperature~0.

\textbf{H0, distributional ceiling.} Neither the minimal prompt nor the expert prompt brings the LLM to the error of the best tabular reference on the four-mode split. Annex~H reports H0 as a paired estimate on the same personas, no equivalence margin having been written before the measurement: the sign, the amplitude and the interval of the gap between each LLM condition and each tabular reference, by cluster resampling at the person level.
The expert-prompt figure is measured once on the sealed cohort, on individuals that were not used to adjust the prompt.

\paragraph{C3. Two regimes no survey tabulates, evaluated for robustness.}

\begin{itemize}
  \item \emph{Five-day hysteresis after a service breakdown.} On day~2 of a five-day horizon, metro line~A fails at 17:00 and imposes a 45-minute delay; the network is nominal again from day~3.
  Three conditions are compared: an agent with a long-term memory store in the sense of \citet{park2023generative}, whose memories decay with a time constant expressed in days; the same agent without memory; the tabular reference, structurally memoryless. No survey follows the same travellers through such an event, so no value is a target. What is testable is the order of the conditions on day~3, \todo{xx.y | exp\_04a vs exp\_04b vs exp\_04c}, and the monotonic recovery over days 3 to 5. The hypothesis falls if the memoryless agent shows the same inertia or if recovery is not monotonic.
  \item \emph{Five real, dated local news events.} Five real, dated articles of the Toulouse local press, chosen before any model call, each act through a channel that no tabular variable represents:
  gusts of the Autan wind above 80\,km/h closing the city's parks; a bed-bug scare in the metro and buses; an open-ended strike of waste collectors in the historic centre; the street parade of La Machine with a pedestrianised centre; the launch of electric bike-sharing on the hilly districts. Each event is run under five conditions on the same personas: the nominal day, with no article; the raw article; a mode-neutral paraphrase, the same fact rewritten without any mention of a mode; a neutral control text, a real local article of similar length with no plausible link to mode choice, which measures the effect of adding any text at all; and the tabular reference with the event encoded into its variables where the event admits such an encoding. For each event and each mode, the expected direction of the shift was written down before any model call, 20 signed predictions in all; we report the share of observed modal shifts that go in the expected direction \todo{xx.y | exp\_05}, with option order randomised for every request, a control made necessary by the order sensitivity SILICA measures \citep{bintareaf2026silica}. The raw-minus-paraphrase contrast isolates instruction following from reasoning; the contrast with the neutral control text separates the effect of the article's content from the effect of adding a text.
\end{itemize}

\subsection{Paper organisation}

\begin{itemize}
  \item Section~2 reviews the related work, from random utility models to generative agents in mobility simulation and to the distributional alignment of LLM populations.
  \item Section~3 describes the device under evaluation: the simulation substrate and the itinerary supply, what the agent sees when it decides, the memory registers, and the chain constraints that keep a simulated day physically consistent.
  \item Section~4 defines the macro and micro metrics, the three rules of the protocol, and the demographic control of the sealed cohort.
  \item Section~5 evaluates behaviour under the minimal prompt and describes the engineering of the expert prompt.
  \item Section~6 confronts the thirteen decision-makers with the survey, measures what prompt engineering moves and on which model, and goes down to the individual decision; the paired differences on which H0 rests are in Annex~H.
  \item Section~7 first sets out the mechanism by which an event reaches a decision, then measures it on two regimes the survey does not tabulate: a shock suffered by one agent, and a press article read by one member of a household and told to the others. The directional predictions are pre-registered.
  \item Section~8 discusses the limits, including the independence-of-irrelevant-alternatives assumption of the renormalisation and the exposure asymmetry between the tabular references and the agents, and the implications for hybrid designs.
  \item Section~9 concludes.
\end{itemize}
